\hypertarget{tuples}{%
\section{Tuples}\label{tuples}}

\hypertarget{un-tuple-est-immuable}{%
\subsection{Un tuple est immuable}\label{un-tuple-est-immuable}}

Un \textbf{tuple} est une séquence de valeurs qui, tout comme une liste,
peut contenir des éléments de n\textquotesingle importe quel type,
indexés par des entiers. La grande différence est que les tuples sont
\textbf{immuables} (\emph{immutable}).

Syntaxiquement, un tuple est une liste de valeurs séparées par des
virgules :

\begin{Shaded}
\begin{Highlighting}[]
\OperatorTok{\textgreater{}\textgreater{}\textgreater{}}\NormalTok{ t }\OperatorTok{=} \StringTok{\textquotesingle{}a\textquotesingle{}}\NormalTok{, }\StringTok{\textquotesingle{}b\textquotesingle{}}\NormalTok{, }\StringTok{\textquotesingle{}c\textquotesingle{}}\NormalTok{, }\StringTok{\textquotesingle{}d\textquotesingle{}}\NormalTok{, }\StringTok{\textquotesingle{}e\textquotesingle{}}
\end{Highlighting}
\end{Shaded}

Bien que ce ne soit pas obligatoire, il est courant de placer les tuples
entre parenthèses :

\begin{Shaded}
\begin{Highlighting}[]
\OperatorTok{\textgreater{}\textgreater{}\textgreater{}}\NormalTok{ t }\OperatorTok{=}\NormalTok{ (}\StringTok{\textquotesingle{}a\textquotesingle{}}\NormalTok{, }\StringTok{\textquotesingle{}b\textquotesingle{}}\NormalTok{, }\StringTok{\textquotesingle{}c\textquotesingle{}}\NormalTok{, }\StringTok{\textquotesingle{}d\textquotesingle{}}\NormalTok{, }\StringTok{\textquotesingle{}e\textquotesingle{}}\NormalTok{)}
\end{Highlighting}
\end{Shaded}

Pour créer un tuple avec un seul élément, vous devez inclure une virgule
finale :

\begin{Shaded}
\begin{Highlighting}[]
\OperatorTok{\textgreater{}\textgreater{}\textgreater{}}\NormalTok{ t1 }\OperatorTok{=}\NormalTok{ (}\StringTok{\textquotesingle{}a\textquotesingle{}}\NormalTok{,)}
\OperatorTok{\textgreater{}\textgreater{}\textgreater{}} \BuiltInTok{type}\NormalTok{(t1)}
\OperatorTok{\textless{}}\KeywordTok{class} \StringTok{\textquotesingle{}tuple\textquotesingle{}}\OperatorTok{\textgreater{}}
\end{Highlighting}
\end{Shaded}

Sans la virgule, Python traite
\texttt{(\textquotesingle{}a\textquotesingle{})} comme une simple chaîne
entre parenthèses :

\begin{Shaded}
\begin{Highlighting}[]
\OperatorTok{\textgreater{}\textgreater{}\textgreater{}}\NormalTok{ t2 }\OperatorTok{=}\NormalTok{ (}\StringTok{\textquotesingle{}a\textquotesingle{}}\NormalTok{)}
\OperatorTok{\textgreater{}\textgreater{}\textgreater{}} \BuiltInTok{type}\NormalTok{(t2)}
\OperatorTok{\textless{}}\KeywordTok{class} \StringTok{\textquotesingle{}str\textquotesingle{}}\OperatorTok{\textgreater{}}
\end{Highlighting}
\end{Shaded}

Pour créer un tuple vide, vous utilisez des parenthèses vides :

\begin{Shaded}
\begin{Highlighting}[]
\OperatorTok{\textgreater{}\textgreater{}\textgreater{}}\NormalTok{ t }\OperatorTok{=}\NormalTok{ ()}
\OperatorTok{\textgreater{}\textgreater{}\textgreater{}} \BuiltInTok{type}\NormalTok{(t)}
\OperatorTok{\textless{}}\KeywordTok{class} \StringTok{\textquotesingle{}tuple\textquotesingle{}}\OperatorTok{\textgreater{}}
\end{Highlighting}
\end{Shaded}

Vous pouvez également utiliser la fonction intégrée \texttt{tuple} pour
créer un tuple à partir d\textquotesingle une séquence :

\begin{Shaded}
\begin{Highlighting}[]
\OperatorTok{\textgreater{}\textgreater{}\textgreater{}}\NormalTok{ t }\OperatorTok{=} \BuiltInTok{tuple}\NormalTok{(}\StringTok{\textquotesingle{}lump\textquotesingle{}}\NormalTok{)}
\OperatorTok{\textgreater{}\textgreater{}\textgreater{}}\NormalTok{ t}
\NormalTok{(}\StringTok{\textquotesingle{}l\textquotesingle{}}\NormalTok{, }\StringTok{\textquotesingle{}u\textquotesingle{}}\NormalTok{, }\StringTok{\textquotesingle{}m\textquotesingle{}}\NormalTok{, }\StringTok{\textquotesingle{}p\textquotesingle{}}\NormalTok{)}
\end{Highlighting}
\end{Shaded}

La plupart des opérateurs de listes fonctionnent également sur les
tuples. L\textquotesingle opérateur de crochet indexe un élément, et
l\textquotesingle opérateur de tranche extrait une section :

\begin{Shaded}
\begin{Highlighting}[]
\OperatorTok{\textgreater{}\textgreater{}\textgreater{}}\NormalTok{ t }\OperatorTok{=}\NormalTok{ (}\StringTok{\textquotesingle{}a\textquotesingle{}}\NormalTok{, }\StringTok{\textquotesingle{}b\textquotesingle{}}\NormalTok{, }\StringTok{\textquotesingle{}c\textquotesingle{}}\NormalTok{, }\StringTok{\textquotesingle{}d\textquotesingle{}}\NormalTok{, }\StringTok{\textquotesingle{}e\textquotesingle{}}\NormalTok{)}
\OperatorTok{\textgreater{}\textgreater{}\textgreater{}}\NormalTok{ t[}\DecValTok{0}\NormalTok{]}
\CommentTok{\textquotesingle{}a\textquotesingle{}}
\OperatorTok{\textgreater{}\textgreater{}\textgreater{}}\NormalTok{ t[}\DecValTok{1}\NormalTok{:}\DecValTok{3}\NormalTok{]}
\NormalTok{(}\StringTok{\textquotesingle{}b\textquotesingle{}}\NormalTok{, }\StringTok{\textquotesingle{}c\textquotesingle{}}\NormalTok{)}
\end{Highlighting}
\end{Shaded}

Cependant, si vous essayez de modifier l\textquotesingle un des éléments
du tuple, vous obtiendrez une erreur, car les tuples sont immuables :

\begin{Shaded}
\begin{Highlighting}[]
\OperatorTok{\textgreater{}\textgreater{}\textgreater{}}\NormalTok{ t[}\DecValTok{0}\NormalTok{] }\OperatorTok{=} \StringTok{\textquotesingle{}A\textquotesingle{}}
\PreprocessorTok{TypeError}\NormalTok{: }\StringTok{\textquotesingle{}tuple\textquotesingle{}} \BuiltInTok{object}\NormalTok{ does }\KeywordTok{not}\NormalTok{ support item assignment}
\end{Highlighting}
\end{Shaded}

Bien que vous ne puissiez pas modifier les éléments d\textquotesingle un
tuple, vous pouvez remplacer un tuple par un autre :

\begin{Shaded}
\begin{Highlighting}[]
\OperatorTok{\textgreater{}\textgreater{}\textgreater{}}\NormalTok{ t }\OperatorTok{=}\NormalTok{ (}\StringTok{\textquotesingle{}A\textquotesingle{}}\NormalTok{,) }\OperatorTok{+}\NormalTok{ t[}\DecValTok{1}\NormalTok{:]}
\OperatorTok{\textgreater{}\textgreater{}\textgreater{}}\NormalTok{ t}
\NormalTok{(}\StringTok{\textquotesingle{}A\textquotesingle{}}\NormalTok{, }\StringTok{\textquotesingle{}b\textquotesingle{}}\NormalTok{, }\StringTok{\textquotesingle{}c\textquotesingle{}}\NormalTok{, }\StringTok{\textquotesingle{}d\textquotesingle{}}\NormalTok{, }\StringTok{\textquotesingle{}e\textquotesingle{}}\NormalTok{)}
\end{Highlighting}
\end{Shaded}

\hypertarget{assignation-de-tuples}{%
\subsection{Assignation de tuples}\label{assignation-de-tuples}}

Il est souvent utile de permuter les valeurs de deux variables. Dans les
langages traditionnels, vous devez utiliser une variable temporaire. En
Python, l\textquotesingle affectation de tuples permet une syntaxe
élégante :

\begin{Shaded}
\begin{Highlighting}[]
\OperatorTok{\textgreater{}\textgreater{}\textgreater{}}\NormalTok{ x, y }\OperatorTok{=}\NormalTok{ y, x}
\end{Highlighting}
\end{Shaded}

Le côté gauche est une séquence de variables ; le côté droit est une
séquence de expressions. Chaque valeur du côté droit est assignée à la
variable correspondante du côté gauche. Toutes les expressions du côté
droit sont évaluées avant toute affectation.

Le nombre de variables à gauche et le nombre de valeurs à droite doivent
être égaux :

\begin{Shaded}
\begin{Highlighting}[]
\OperatorTok{\textgreater{}\textgreater{}\textgreater{}}\NormalTok{ x, y }\OperatorTok{=} \DecValTok{1}\NormalTok{, }\DecValTok{2}\NormalTok{, }\DecValTok{3}
\PreprocessorTok{ValueError}\NormalTok{: too many values to unpack}
\end{Highlighting}
\end{Shaded}

Plus généralement, l\textquotesingle affectation de tuples fonctionne
avec n\textquotesingle importe quelle séquence.

\hypertarget{les-tuples-comme-valeurs-de-retour}{%
\subsection{Les tuples comme valeurs de
retour}\label{les-tuples-comme-valeurs-de-retour}}

Une fonction ne peut retourner qu\textquotesingle une seule valeur, mais
si cette valeur est un tuple, l\textquotesingle effet est le même que de
retourner plusieurs valeurs. Par exemple, la fonction \texttt{divmod}
prend deux entiers et renvoie un tuple de deux valeurs : le quotient et
le reste.

\begin{Shaded}
\begin{Highlighting}[]
\OperatorTok{\textgreater{}\textgreater{}\textgreater{}}\NormalTok{ t }\OperatorTok{=} \BuiltInTok{divmod}\NormalTok{(}\DecValTok{7}\NormalTok{, }\DecValTok{3}\NormalTok{)}
\OperatorTok{\textgreater{}\textgreater{}\textgreater{}}\NormalTok{ t}
\NormalTok{(}\DecValTok{2}\NormalTok{, }\DecValTok{1}\NormalTok{)}
\end{Highlighting}
\end{Shaded}

Vous pouvez également utiliser l\textquotesingle affectation de tuples
pour stocker les éléments séparément :

\begin{Shaded}
\begin{Highlighting}[]
\OperatorTok{\textgreater{}\textgreater{}\textgreater{}}\NormalTok{ quotient, reste }\OperatorTok{=} \BuiltInTok{divmod}\NormalTok{(}\DecValTok{7}\NormalTok{, }\DecValTok{3}\NormalTok{)}
\OperatorTok{\textgreater{}\textgreater{}\textgreater{}}\NormalTok{ quotient}
\DecValTok{2}
\OperatorTok{\textgreater{}\textgreater{}\textgreater{}}\NormalTok{ reste}
\DecValTok{1}
\end{Highlighting}
\end{Shaded}

Un autre exemple est la fonction \texttt{min\_max}, qui trouve à la fois
la plus petite et la plus grande valeur d\textquotesingle une séquence :

\begin{Shaded}
\begin{Highlighting}[]
\KeywordTok{def}\NormalTok{ min\_max(t):}
    \ControlFlowTok{return} \BuiltInTok{min}\NormalTok{(t), }\BuiltInTok{max}\NormalTok{(t)}
\end{Highlighting}
\end{Shaded}

\hypertarget{tuples-uxe0-longueur-variable-darguments}{%
\subsection{Tuples à longueur variable
d\textquotesingle arguments}\label{tuples-uxe0-longueur-variable-darguments}}

Les fonctions peuvent prendre un nombre variable
d\textquotesingle arguments. Un paramètre commençant par un astérisque
(\texttt{*}) \textbf{rassemble} les arguments en un tuple. Par exemple,
\texttt{somme\_tous} prend un nombre quelconque
d\textquotesingle arguments et calcule leur somme :

\begin{Shaded}
\begin{Highlighting}[]
\KeywordTok{def}\NormalTok{ somme\_tous(}\OperatorTok{*}\NormalTok{args):}
    \ControlFlowTok{return} \BuiltInTok{sum}\NormalTok{(args)}

\OperatorTok{\textgreater{}\textgreater{}\textgreater{}}\NormalTok{ somme\_tous(}\DecValTok{1}\NormalTok{, }\DecValTok{2}\NormalTok{, }\DecValTok{3}\NormalTok{)}
\DecValTok{6}
\end{Highlighting}
\end{Shaded}

À l\textquotesingle inverse, l\textquotesingle opérateur \texttt{*} peut
être utilisé pour \textbf{répartir} une séquence en tant
qu\textquotesingle arguments d\textquotesingle une fonction :

\begin{Shaded}
\begin{Highlighting}[]
\OperatorTok{\textgreater{}\textgreater{}\textgreater{}}\NormalTok{ t }\OperatorTok{=}\NormalTok{ (}\DecValTok{7}\NormalTok{, }\DecValTok{3}\NormalTok{)}
\OperatorTok{\textgreater{}\textgreater{}\textgreater{}} \BuiltInTok{divmod}\NormalTok{(}\OperatorTok{*}\NormalTok{t)}
\NormalTok{(}\DecValTok{2}\NormalTok{, }\DecValTok{1}\NormalTok{)}
\end{Highlighting}
\end{Shaded}

\hypertarget{listes-et-tuples}{%
\subsection{Listes et tuples}\label{listes-et-tuples}}

La fonction \texttt{zip} prend deux ou plusieurs séquences et les
combine en un itérateur de tuples, où chaque tuple contient un élément
de chaque séquence.

\begin{Shaded}
\begin{Highlighting}[]
\OperatorTok{\textgreater{}\textgreater{}\textgreater{}}\NormalTok{ s }\OperatorTok{=} \StringTok{\textquotesingle{}abc\textquotesingle{}}
\OperatorTok{\textgreater{}\textgreater{}\textgreater{}}\NormalTok{ t }\OperatorTok{=}\NormalTok{ [}\DecValTok{1}\NormalTok{, }\DecValTok{2}\NormalTok{, }\DecValTok{3}\NormalTok{]}
\OperatorTok{\textgreater{}\textgreater{}\textgreater{}} \BuiltInTok{zip}\NormalTok{(s, t)}
\OperatorTok{\textless{}}\BuiltInTok{zip} \BuiltInTok{object}\NormalTok{ at }\DecValTok{0}\ErrorTok{x}\NormalTok{...}\OperatorTok{\textgreater{}}
\end{Highlighting}
\end{Shaded}

Pour utiliser le résultat de \texttt{zip}, vous pouvez le convertir en
liste de tuples ou utiliser une boucle \texttt{for} :

\begin{Shaded}
\begin{Highlighting}[]
\OperatorTok{\textgreater{}\textgreater{}\textgreater{}} \BuiltInTok{list}\NormalTok{(}\BuiltInTok{zip}\NormalTok{(s, t))}
\NormalTok{[(}\StringTok{\textquotesingle{}a\textquotesingle{}}\NormalTok{, }\DecValTok{1}\NormalTok{), (}\StringTok{\textquotesingle{}b\textquotesingle{}}\NormalTok{, }\DecValTok{2}\NormalTok{), (}\StringTok{\textquotesingle{}c\textquotesingle{}}\NormalTok{, }\DecValTok{3}\NormalTok{)]}
\end{Highlighting}
\end{Shaded}

Si les longueurs des séquences sont différentes, le résultat
s\textquotesingle arrête à la longueur de la plus courte.

Vous pouvez parcourir les éléments d\textquotesingle une liste de tuples
en utilisant une affectation de tuples dans une boucle \texttt{for} :

\begin{Shaded}
\begin{Highlighting}[]
\NormalTok{t }\OperatorTok{=}\NormalTok{ [(}\StringTok{\textquotesingle{}a\textquotesingle{}}\NormalTok{, }\DecValTok{1}\NormalTok{), (}\StringTok{\textquotesingle{}b\textquotesingle{}}\NormalTok{, }\DecValTok{2}\NormalTok{), (}\StringTok{\textquotesingle{}c\textquotesingle{}}\NormalTok{, }\DecValTok{3}\NormalTok{)]}
\ControlFlowTok{for}\NormalTok{ lettre, nombre }\KeywordTok{in}\NormalTok{ t:}
    \BuiltInTok{print}\NormalTok{(nombre, lettre)}
\end{Highlighting}
\end{Shaded}

\hypertarget{dictionnaires-et-tuples}{%
\subsection{Dictionnaires et tuples}\label{dictionnaires-et-tuples}}

Les dictionnaires possèdent une méthode appelée \texttt{items} qui
renvoie une séquence de paires clé-valeur sous forme de tuples.

\begin{Shaded}
\begin{Highlighting}[]
\OperatorTok{\textgreater{}\textgreater{}\textgreater{}}\NormalTok{ d }\OperatorTok{=}\NormalTok{ \{}\StringTok{\textquotesingle{}a\textquotesingle{}}\NormalTok{: }\DecValTok{0}\NormalTok{, }\StringTok{\textquotesingle{}b\textquotesingle{}}\NormalTok{: }\DecValTok{1}\NormalTok{, }\StringTok{\textquotesingle{}c\textquotesingle{}}\NormalTok{: }\DecValTok{2}\NormalTok{\}}
\OperatorTok{\textgreater{}\textgreater{}\textgreater{}}\NormalTok{ t }\OperatorTok{=} \BuiltInTok{list}\NormalTok{(d.items())}
\OperatorTok{\textgreater{}\textgreater{}\textgreater{}}\NormalTok{ t}
\NormalTok{[(}\StringTok{\textquotesingle{}a\textquotesingle{}}\NormalTok{, }\DecValTok{0}\NormalTok{), (}\StringTok{\textquotesingle{}c\textquotesingle{}}\NormalTok{, }\DecValTok{2}\NormalTok{), (}\StringTok{\textquotesingle{}b\textquotesingle{}}\NormalTok{, }\DecValTok{1}\NormalTok{)]}
\end{Highlighting}
\end{Shaded}

À l\textquotesingle inverse, vous pouvez utiliser une liste de tuples
pour initialiser un nouveau dictionnaire :

\begin{Shaded}
\begin{Highlighting}[]
\OperatorTok{\textgreater{}\textgreater{}\textgreater{}}\NormalTok{ d }\OperatorTok{=} \BuiltInTok{dict}\NormalTok{([(}\StringTok{\textquotesingle{}a\textquotesingle{}}\NormalTok{, }\DecValTok{0}\NormalTok{), (}\StringTok{\textquotesingle{}b\textquotesingle{}}\NormalTok{, }\DecValTok{1}\NormalTok{), (}\StringTok{\textquotesingle{}c\textquotesingle{}}\NormalTok{, }\DecValTok{2}\NormalTok{)])}
\OperatorTok{\textgreater{}\textgreater{}\textgreater{}}\NormalTok{ d}
\NormalTok{\{}\StringTok{\textquotesingle{}a\textquotesingle{}}\NormalTok{: }\DecValTok{0}\NormalTok{, }\StringTok{\textquotesingle{}b\textquotesingle{}}\NormalTok{: }\DecValTok{1}\NormalTok{, }\StringTok{\textquotesingle{}c\textquotesingle{}}\NormalTok{: }\DecValTok{2}\NormalTok{\}}
\end{Highlighting}
\end{Shaded}

La combinaison de \texttt{dict} et \texttt{zip} permet de créer
rapidement des dictionnaires :

\begin{Shaded}
\begin{Highlighting}[]
\OperatorTok{\textgreater{}\textgreater{}\textgreater{}}\NormalTok{ d }\OperatorTok{=} \BuiltInTok{dict}\NormalTok{(}\BuiltInTok{zip}\NormalTok{(}\StringTok{\textquotesingle{}abc\textquotesingle{}}\NormalTok{, }\BuiltInTok{range}\NormalTok{(}\DecValTok{3}\NormalTok{)))}
\OperatorTok{\textgreater{}\textgreater{}\textgreater{}}\NormalTok{ d}
\NormalTok{\{}\StringTok{\textquotesingle{}a\textquotesingle{}}\NormalTok{: }\DecValTok{0}\NormalTok{, }\StringTok{\textquotesingle{}b\textquotesingle{}}\NormalTok{: }\DecValTok{1}\NormalTok{, }\StringTok{\textquotesingle{}c\textquotesingle{}}\NormalTok{: }\DecValTok{2}\NormalTok{\}}
\end{Highlighting}
\end{Shaded}

\hypertarget{glossaire}{%
\subsection{Glossaire}\label{glossaire}}

tuple tuple Une séquence immuable d\textquotesingle éléments.

affectation de tuples tuple assignment Une affectation qui associe un
tuple de valeurs à une séquence de variables.

rassembler gather Le processus de regroupement
d\textquotesingle arguments à longueur variable dans un tuple.

répartir scatter Le processus de traitement d\textquotesingle une
séquence de valeurs en tant qu\textquotesingle arguments de fonction.

\hypertarget{exercices}{%
\subsection{Exercices}\label{exercices}}

\textbf{Exercice 1}

Écrivez une fonction appelée \texttt{somme\_tous} qui prend un nombre
quelconque de nombres en arguments et renvoie leur somme.

\textbf{Exercice 2}

Modifiez un programme de tri de mots de sorte qu\textquotesingle il trie
les mots par longueur, du plus long au plus court, et que les mots de
même longueur soient triés par ordre alphabétique.

\textbf{Exercice 3}

Écrivez une fonction appelée \texttt{frequence\_lettres} qui prend une
chaîne de caractères et renvoie un dictionnaire des lettres comptées,
triées par ordre décroissant de fréquence.

\textbf{Exercice 4}

Un anagramme est un mot formé en réorganisant les lettres
d\textquotesingle un autre mot. Écrivez un programme qui lit une liste
de mots à partir d\textquotesingle un fichier et imprime tous les
ensembles de mots qui sont des anagrammes entre eux.
